\subsection{I 智能体编码评测}\label{i-agentic-coding-evaluations}

SWE-bench Verified（Chowdhury 等，2024）、SWE-Bench Pro（Deng 等，2025）和
Terminal-Bench 2.0（Merrill 等，2026）的智能体编码基准结果见表 11 的第二部分。
与上一节的 STEM 评测不同，智能体编码评测是多轮交互式评测，要求模型与环境进行
交互。对于这三个基准，我们都使用一种非常简单的、ReAct 风格（Yao 等，2023）的
"始终追加"（always append）智能体循环来评测我们的模型，该循环描述见图 18。
在每一轮中，模型可以生成工具调用，我们会在环境中执行这些调用，并在模型发起下
一次工具调用之前将输出追加到轨迹中。评测时我们使用 256k token 的总上下文长度；
SWE-bench Verified 和 SWE-Bench Pro 的最大输出长度为 8k token，Terminal-Bench
2.0 为 32k token，最大步数为 1,000 步。当模型不再输出工具调用、达到上下文长度
上限或达到步数上限时，智能体循环即告终止。循环终止后，相应的评分器会在智能体
所使用的 SEE 环境中运行。对于 SWE-bench Verified 和 SWE-Bench Pro，我们使用与
爬山过程中相同的两种工具（bash 和 string replace editor）进行评测，如第 3.3.1
节所述。对于 Terminal-Bench 2.0，我们仅使用 bash 工具进行评测。

SWE-bench Verified。SWE-bench Verified 是 SWE-Bench（Jimenez 等，2024）中人工
精选的 500 个任务子集，旨在提高评测的可靠性。每个任务都来自热门 Python 仓库中
的真实公开 GitHub issue，模型必须生成一个能解决该 issue 并通过仓库隐藏测试套件
的补丁。所有任务均经过审核，以确保问题陈述可解且测试能公正地评估正确性。

SWE-Bench Pro。SWE-Bench Pro 是 SWE-bench Verified 更具挑战性的继任者，包含
731 个任务，采用相同的基于单元测试的评分机制，旨在评估 LLM 在来自真实世界开源
仓库的专业级软件工程任务上的表现。它强调更困难的多文件修改、更长时域的推理，
以及更贴近工程师实际工作的复杂 bug 修复或功能实现。更重要的是，与 SWE-bench
Verified 不同，SWE-Bench Pro 包含 Python 之外其他语言的仓库。

Terminal-Bench 2.0。Terminal-Bench 2.0 是 Terminal-Bench（Merrill 等，2026）
经人工验证、难度大幅提升的继任者，包含覆盖软件工程、调试、数据科学、机器学习、
安全和系统管理等领域的 89 个真实任务。每个任务都在沙箱化终端环境中运行，并根据
输出或环境最终状态进行程序化评分。任务设计旨在考查长期规划、错误恢复和真正的
命令行熟练度，而非模式匹配。

\subsection{J 安全评测}\label{j-safety-evaluations}

\subsection{J.1 内部评测细节}\label{j.1-internal-evaluation-details}

\subsection{J.1.1 方法论}\label{j.1.1-methodology}

评测分两个阶段进行：先对用户请求进行结构化刻画，再据此推导出用于给模型行为评分
的响应规范。这种分解使我们能够区分真正的安全缺陷与在良性但敏感话题上的过度谨慎
行为，并给出按策略、按模态的细分指标，而非单一的汇总分数。

表 18. 用户请求分类维度。

{\def\LTcaptype{none} % do not increment counter
\begin{longtable}[]{@{}|l|l|@{}}
\toprule\noalign{}
\endhead
\bottomrule\noalign{}
\endlastfoot
\hline
维度 & 描述 \\
\hline
危害类别 \& 策略违规 & 识别请求涉及哪个策略领域，以及是否违反特定规则，
将判断建立在具体策略条款之上。 \\
\hline
敏感度等级 & 基于影响范围（个体 vs. 广泛）和严重程度（最坏情况下的滥用后果），
按序数尺度记录主题敏感度，与请求是否被禁止无关。 \\
\hline
用户意图 & 二值信号：可能存在善意意图 vs. 恶意意图。仅话题危险本身不构成
恶意的证据。 \\
\hline
用户场景 & 捕捉用户的实用目标（例如：信息检索、寻求指导、创意生成、
情景扮演）。 \\
\hline
越狱检测 & 双重作用：筛查对抗性包装，并为请求是否以对抗方式构造提供信号。 \\
\hline
\end{longtable}
}

两个评测阶段均由 LLM 评判器执行，其中分类步骤会运行多次并按多数投票聚合，以
控制评判器的方差。每个用户请求还会经过越狱检测，该检测利用攻击族与技术变体的
分层分类体系筛查对抗性包装。一旦检测到包装，提取步骤会还原底层的原始请求，
后续分类基于还原后的请求进行。这确保评测反映的是模型对实际底层意图的行为，而非
对表层攻击模板的行为，从而防止针对特定包装的过拟合虚增安全分数。

阶段 1：用户请求分类。每个用户请求都按照表 18 所述的五种互补维度进行刻画。这些
维度共同捕捉请求中与策略相关的内容以及请求所处的上下文。对于多轮对话，整个对话
历史都被视为上下文输入。

阶段 2：响应策略生成。阶段 1 的分类结果由评分规则评判器（rubric judge）消费，
该评判器针对每个请求输出响应规范，涵盖相关的策略约束、适当的拒答姿态（根据意图
和严重程度，可以是完全遵从、部分拒答或完全拒答），以及响应所要求的语气和方式。
对于被归类为无害或意图不明的请求，该 LLM 生成的评分规则会被一种确定性的默认
规则取代，后者会明确惩罚含糊其辞和无端的拒答。

\subsection{J.1.2 数据收集与采样}\label{j.1.2-data-collection-sampling}

设计原则：评测结果的信息量取决于用于评测的提示词本身：如果评测集过于狭窄、
权重失当或已经饱和，就会围绕错误的量产生狭窄的置信区间，任何基于该评测集的模型
选择决策都会继承这一错误。我们将评测集的构建视为一个方法论问题，评测集必须满足
三个性质（Atkinson 等，2007）。

• 覆盖上文介绍的政策领域、敏感度等级、用户场景和对抗性框架，使任何行为模式不同
的输入切片都不会在无声无息中代表性不足。

• 在固定评测预算下的统计效率，因为评判成本高昂，且各类别指标的置信区间随样本量
增大的收窄速度很慢。

• 与其他前沿模型比较时的区分能力，使模型改进时总分确实会变化，而不会因提示词
过于简单或过于困难而饱和。

来源池与标注：候选提示词池由三个互补来源汇集而成：公开安全基准、已过滤 PII
（个人身份信息）的消费者 Copilot 日志，以及外部供应商根据我们的策略分类体系
收集的提示词。每个候选提示词都会经过上文所述的标注框架，为后续抽样所需的 5 个
分层附加标签：危害类别、策略违规、敏感度、用户场景和越狱。

信息量加权实验设计。我们力求从候选提示词池中选出信息量最大的提示词用于评测。
我们将此视为一个自动化测试组装问题来处理。这与近期通过心理测量学建模大幅压缩
LLM 评测成本的努力方向一致，例如 tinyBenchmarks（Polo 等，2024），它利用项目
反应理论（Item Response Theory）来识别静态的、具有普遍代表性的评测子集。作为
该框架的扩展，我们对模型并非始终失败或始终通过的提示词赋予更高权重。这有助于
我们主动发现最有可能改进模型的方向。

阈值设定。上文所述的评测集是安全考量进入 MAI-Thinking-1 爬山过程的主要机制。
每个内部发布候选版本都必须在一组固定的安全评测上达到验收阈值，并且不得相对上一
版已验收的发布版本发生回退。这组评测既涵盖针对性的策略领域表现，也涵盖在无偏
排行榜集上的整体行为，从而避免仅针对单一危害类别的回退被其他方面的改进平均掉。
阈值在流程的两个节点发挥作用：其一，作为发布候选签核（sign-off）时的门槛，
任何在组评测指标上低于阈值的候选版本都无资格进入发布配方；其二，当实验模型低于
较低阈值时触发预警警报。

阈值本身源自对安全、过度拒答和下游模型质量之间可达权衡曲面的帕累托分析：在每个
发布周期，我们绘制该周期内所评估配方的前沿面，并识别出拐点，在拐点处边际安全
收益开始以不成比例的过度拒答或质量损失为代价。随后，阈值被设定为沿该前沿面的
固定百分位数，而非静态的绝对数值，这样门槛就能跟随当前架构与数据配比下可实现
的水平，而不是一个过时的历史下限。这使得门槛信号在不同发布周期之间的含义保持
稳定 --- 相对于当前前沿面，通过门槛即提供同样的安全质量证据，无论前沿面如何
移动 --- 同时为绝对数值随着底层配方的改进而上升留有空间。
